\section{Results and interpretation}
The requested target is a fixed ordinary integer polynomial $P$, on a stated domain $D$, for which
\[
 \ZFC\vdash \bigl(\neg\exists\mathbf x\in D^n\ P(\mathbf x)=0\bigr)
 \longleftrightarrow\Con(\ZFC).
\]
Here consistency is syntactic: there is no finite derivation of a contradiction in the fixed effective proof calculus for ZFC. We prove an existence equivalence to such a refutation, then negate it. No argument assumes that ZFC is consistent, and no part of the work decides that question.

\begin{center}
\begin{tabular}{llrr}
\toprule
Domain & New construction & Operations & Exact degree\\
\midrule
$\Zp^9$ & Gap $N-q$, simple majorant & 444 & 1,845,940\\
$\Zp^9$ & Gap $N$, simple majorant & 449 & 1,834,516\\
$\Zp^9$ & Reflected gap, simple majorant & 451 & 1,825,540\\
$\Zp^9$ & Gap $N-q$, refined majorant & 476 & 1,520,178\\
$\Zp^9$ & Gap $N$, refined majorant & 481 & 1,510,770\\
$\Zp^9$ & Reflected gap, refined majorant & 483 & 1,503,378\\
$\Z^{10}$ & Reflected gap, congruence norm & 494 & 3,507,876\\
$\Z^{10}$ & Gap $N-q$, congruence norm & 490 & 3,547,076\\
$\Zp^{15}$ & Integral-root transport & 88 & 24\\
$\Zp^{15}$ & Integral-root transport & 87 & 26\\
$\Zp^{15}$ & Integral-root identity & 78 & 30\\
\bottomrule
\end{tabular}
\end{center}
The nine-variable family additionally has explicit nonnegative-domain versions, each with seven extra operations and the same degree. Against the newest incoming nine-positive circuit (473 operations, degree 1,910,404), the 444-operation variant improves both objectives; the refined variants trade some gates for lower degree. Against the incoming ten-integer degree/operation records considered here, both completed ten-integer variants improve substantially. The ten-integer evaluator materializes a fixed threshold constant: retaining the former five constant-only gates would give 499 and 495 operations respectively. This convention does not change either polynomial or its degree.


Every named intermediate register in the completed equations is a polynomial abbreviation, not another quantified variable. Every exponent is a fixed nonnegative integer. There are no divisions, variable powers, square roots, bit operations or unbounded quantifiers inside the final literal circuits. Divisions, radicals and population counts occur only in proofs. The four compiler ports are bound to complete, coherent numerical tuples. Derived large coefficients are materialized as signed binary payloads with hashes rather than left as unspecified universal-polynomial parameters.

\paragraph{What simplicity means here.}
Number of unknowns, their domain, total polynomial degree, arithmetic circuit cost, and coefficient storage are different objectives. We count a live binary addition, subtraction or multiplication as one operation; multiplication by a numeral is charged, and fixed powers are compiled to multiplication gates. Reusing a previously computed node is free. Fixed integer coefficient construction is not included in this circuit cost. Thus a smaller operation count is not a bit-complexity or running-time claim. We retain the relevant tradeoffs rather than declare a single globally simplest polynomial.

\paragraph{A route beyond the previous architecture.}
The new direct proof source does not use the K215 transition table or the Jones universal expression. It addresses the user's broader permission to choose another equiconsistent formalism. Its local transition polynomial and accepting-path equivalence are explicit. The report does not misidentify that local polynomial with a completed consistency equation: converting all finite paths to one fixed ordinary polynomial still has a separately stated implementation boundary.
